% 第三章：整体框架
% 负责人：马福泉

\section{整体框架} \label{sec:design}

\subsection{系统架构}

Yat-CASched调度器采用分层架构设计，主要包括以下模块：

\begin{itemize}
    \item \textbf{任务管理模块}：负责DAG任务的解析、存储和管理
    \item \textbf{缓存建模模块}：实现Cache Recency Profile (CRP)预测模型
    \item \textbf{调度决策模块}：基于MSF和LCIF规则进行核心分配
    \item \textbf{运行时监控模块}：跟踪任务执行状态和缓存使用情况
    \item \textbf{性能评估模块}：收集性能指标并进行分析
\end{itemize}

\subsection{数据结构设计}

\subsubsection{任务表示}

\begin{code}
struct DAGTask {
    int task_id;
    vector<Node> nodes;
    vector<Edge> edges;
    int period;
    int priority;
    int total_workload;
};

struct Node {
    int node_id;
    int wcet;
    vector<int> predecessors;
    vector<int> successors;
};
\end{code}

\subsubsection{调度状态}

\begin{code}
struct SchedulingState {
    vector<Core> cores;
    map<int, int> allocation_history;
    SpeedupTable speedup_table;
    CacheRecencyProfile crp_model;
};
\end{code}

\subsection{接口设计}

调度器提供以下主要接口：

\begin{itemize}
    \item \texttt{schedule\_tasks()}: 主调度函数
    \item \texttt{update\_crp\_model()}: 更新CRP模型
    \item \texttt{calculate\_speedup()}: 计算加速比
    \item \texttt{allocate\_core()}: 核心分配决策
\end{itemize}
